\documentclass{article}

\usepackage{bm}
\usepackage{float}
\usepackage{xcolor}
\usepackage{booktabs}
\usepackage{setspace}
\usepackage{datetime}
\usepackage{geometry}
\usepackage{graphicx}
\usepackage{mathrsfs}
\usepackage{listings}
\usepackage{tcolorbox}
\usepackage{amsmath, amssymb, amsfonts}
\usepackage[fontsize = 12pt]{fontsize}

\newcommand{\mathsize}{\fontsize{7pt}{9pt}\selectfont}
\renewcommand{\thesection}{\Roman{section}}
\renewcommand{\thesubsection}{\arabic{subsection}}
\renewcommand{\thesubsubsection}{(\alph{subsubsection})}
\newcommand{\diff}[2]{\dfrac{\mathrm{d}#1}{\mathrm{d}#2}}
\newcommand{\diffn}[3]{\dfrac{\mathrm{d}^{#3}#1}{\mathrm{d}#2^{#3}}}
\newcommand{\piff}[2]{\dfrac{\partial#1}{\partial#2}}
\newcommand{\piffn}[3]{\dfrac{\partial^{#3}#1}{\partial#2^{#3}}}
\newcommand{\eval}[2]{\left.#1\right|_{#2}}
\newcommand{\e}{\mathrm{e}}
\renewcommand{\d}{\mathrm{d}}

\definecolor{gray}{RGB}{128, 128, 128}

\setcounter{MaxMatrixCols}{20}

\geometry{
    a4paper,
    left = 1.75cm,
    right = 1.75cm,
    top = 2cm,
    bottom = 2cm
}

\title{Data Mining: Homework 3}
\author{Illusionna \quad 2025..............4 \quad Artificial Intelligence School}
\date{\currenttime,\ \today}



\begin{document}

\maketitle
\begin{spacing}{2}
    \tableofcontents
\end{spacing}
\thispagestyle{empty}
\clearpage
\setcounter{page}{1}


\section{Written Assignment}
\subsection{K-means}
\textcolor{gray}{\it Suppose that the data mining task is to cluster the following ten points (with(x, y, z) representing location) into three clusters:}
\textcolor{gray}{
    \begin{align*}
        \rm\color{gray}& A1(4,2,5) \quad A2(10,5,2) \quad A3(5,8,7)
        \\
        \rm& B1(1,1,1) \quad B2(2,3,2) \quad B3(3,6,9)
        \\
        \rm& C1(11,9,2) \quad C2(1,4,6) \quad C3(9,1,7) \quad C4(5,6,7)
    \end{align*}
}
\textcolor{gray}{\it The distance function is Euclidean distance. Suppose initially we assign A2,B2,C2 as the center of each cluster, respectively. Use the K-Means algorithm to show only.}



\subsubsection{Centroids}
\textcolor{gray}{\it The three cluster's centers after the first round execution.}

Define the Euclidean distance as:
\[ d(\bm{x},\ \bm{y})=\sqrt{\displaystyle\sum_{i=1}^n(x_i-y_i)^2}\implies\sqrt{\displaystyle(x_1-y_1)^2+(x_2-y_2)^2+(x_3-y_3)^2} \]
\begin{align*}
    \begin{matrix}
         & \rm A2(10,5,2) & \rm B2(2,3,2) & \rm C2(1,4,6) & \rm Centroid \\
        \rm \color{blue}A1(4,2,5) & 7.3485 & 3.7417 & 3.7417 & \rm \color{blue}B2 \\
        \rm \color{green}A3(5,8,7) & 7.6811 & 7.6811 & 5.7446 & \rm \color{green}C2 \\
        \rm \color{blue}B1(1,1,1) & 9.8995 & 2.4495 & 5.8310 & \rm \color{blue}B2 \\
        \rm \color{green}B3(3,6,9) & 9.9499 & 7.6811 & 4.1231 & \rm \color{green}C2 \\
        \rm \color{red}C1(11,9,2) & 4.1231 & 10.8167 & 11.8743 & \rm \color{red}A2 \\
        \rm \color{red}C3(9,1,7) & 6.4807 & 8.8318 & 8.6023 & \rm \color{red}A2 \\
        \rm \color{green}C4(5,6,7) & 7.1414 & 6.5574 & 4.5826 & \rm \color{green}C2 \\
    \end{matrix}
\end{align*}

The three cluster's centers after the first iteration are \{A1, B1, B2\}, \{A2, C1, C3\} and \{A3, B3, C2, C4\}.



\subsubsection{Clusters}
\textcolor{gray}{\it The final three clusters.}

The latest centroids are K1, K2 and K3.
\[ \rm K1=\dfrac{A1+B1+B2}{3}\approx(2.3333,2.0000,2.6667) \]
\[ \rm K2=\dfrac{A2+C1+C3}{3}\approx(10.0000, 5.0000, 3.6667) \]
\[ \rm K3=\dfrac{A3+B3+C2+C4}{4}=(3.5000,6.0000,7.2500) \]


\begin{align*}
    \begin{matrix}
         & \rm K1 & \rm K2 & \rm K3 & \rm Centroid \\
        \rm \color{red}A1 & 2.8674 & 6.8394 & 4.6165 & \rm \color{red}K1 \\
        \rm \color{green}A2 & 8.2597 & 1.6667 & 8.4150 & \rm \color{green}K2 \\
        \rm \color{blue}A3 & 7.8669 & 6.7165 & 2.5125 & \rm \color{blue}K3 \\
        \rm \color{red}B1 & 2.3570 & 10.2035 & 8.3853 & \rm \color{red}K1 \\
        \rm \color{red}B2 & 1.2472 & 8.4130 & 6.2300 & \rm \color{red}K1 \\
        \rm \color{blue}B3 & 7.5203 & 8.8569 & 1.8200 & \rm \color{blue}K3 \\
        \rm \color{green}C1 & 11.1604 & 4.4472 & 9.6339 & \rm \color{green}K2 \\
        \rm \color{blue}C2 & 4.1096 & 9.3512 & 3.4369 & \rm \color{blue}K3 \\
        \rm \color{green}C3 & 8.0139 & 5.3020 & 7.4372 & \rm \color{green}K2 \\
        \rm \color{blue}C4 & 6.4722 & 6.0919 & 1.5207 & \rm \color{blue}K3 \\
    \end{matrix}
\end{align*}

The latest three cluster's centers after the second iteration are \{A1, B1, B2\}, \{A2, C1, C3\} and \{A3, B3, C2, C4\}. The centers of clusters have remained \textbf{unchanged} and \textbf{converged}.



\subsection{AGNES}
\textcolor{gray}{\it For the data set in Q1, perform AGNES algorithm on the data set to generate 2 clusters. Note: The two clusters with the minimum Single-Link are merged at each step. Please list the clusters to be merged at each step.}

Minimum Single-Link:
\begin{align*}
    \begin{matrix}
         & \{A1\} & \{A2\} & \{A3\} & \{B1\} & \{B2\} & \{B3\} & \{C1\} & \{C2\} & \{C3\} & \{C4\} \\
        \{A1\} & 0.0000 \\
        \{A2\} & 7.3485 & 0.0000 \\
        \{A3\} & 6.4031 & 7.6811 & 0.0000 \\
        \{B1\} & 5.0990 & 9.8995 & 10.0499 & 0.0000 \\
        \{B2\} & 3.7417 & 8.2462 & 7.6811 & 2.4495 & 0.0000 \\
        \{B3\} & 5.7446 & 9.9499 & 3.4641 & 9.6437 & 7.6811 & 0.0000 \\
        \{C1\} & 10.3441 & 4.1231 & 7.8740 & 12.8452 & 10.8167 & 11.0454 & 0.0000 \\
        \{C2\} & 3.7417 & 9.8995 & 5.7446 & 5.8310 & 4.2426 & 4.1231 & 11.8743 & 0.0000 \\
        \{C3\} & 5.4772 & 6.4807 & 8.0623 & 10.0000 & 8.8318 & 8.0623 & 9.6437 & 8.6023 & 0.0000 \\
        \{C4\} & 4.5826 & 7.1414 & $\color{red}2.0000$ & 8.7750 & 6.5574 & 2.8284 & 8.3666 & 4.5826 & 6.4031 & 0.0000 \\
    \end{matrix}
\end{align*}

Merge clusters \{A3\} and \{C4\}.
\begin{align*}
    \begin{matrix}
         & \{A1\} & \{A2\} & \{B1\} & \{B2\} & \{B3\} & \{C1\} & \{C2\} & \{C3\} & \{A3,C4\} \\
        \{A1\} & 0.0000 \\
        \{A2\} & 7.3485 & 0.0000 \\
        \{B1\} & 5.0990 & 9.8995 & 0.0000 \\
        \{B2\} & 3.7417 & 8.2462 & $\color{red}2.4495$ & 0.0000 \\
        \{B3\} & 5.7446 & 9.9499 & 9.6437 & 7.6811 & 0.0000 \\
        \{C1\} & 10.3441 & 4.1231 & 12.8452 & 10.8167 & 11.0454 & 0.0000 \\
        \{C2\} & 3.7417 & 9.8995 & 5.8310 & 4.2426 & 4.1231 & 11.8743 & 0.0000 \\
        \{C3\} & 5.4772 & 6.4807 & 10.0000 & 8.8318 & 8.0623 & 9.6437 & 8.6023 & 0.0000 \\
        \{A3,C4\} & 4.5826 & 7.1414 & 8.7750 & 6.5574 & 2.8284 & 7.8740 & 4.5826 & 6.4031 & 0.0000 \\
    \end{matrix}
\end{align*}

Merge clusters \{B1\} and \{B2\}.

\begin{align*}
    \begin{matrix}
         & \{A1\} & \{A2\} & \{B3\} & \{C1\} & \{C2\} & \{C3\} & \{A3,C4\} & \{B1,B2\} \\
        \{A1\} & 0.0000 \\
        \{A2\} & 7.3485 & 0.0000 \\
        \{B3\} & 5.7446 & 9.9499 & 0.0000 \\
        \{C1\} & 10.3441 & 4.1231 & 11.0454 & 0.0000 \\
        \{C2\} & 3.7417 & 9.8995 & 4.1231 & 11.8743 & 0.0000 \\
        \{C3\} & 5.4772 & 6.4807 & 8.0623 & 9.6437 & 8.6023 & 0.0000 \\
        \{A3,C4\} & 4.5826 & 7.1414 & $\color{red}2.8284$ & 7.8740 & 4.5826 & 6.4031 & 0.0000 \\
        \{B1,B2\} & 3.7417 & 8.2462 & 7.6811 & 10.8167 & 4.2426 & 8.8318 & 6.5574 & 0.0000 \\
    \end{matrix}
\end{align*}

Merge clusters \{B3\} and \{A3, C4\}.
\begin{align*}
    \begin{matrix}
         & \{A1\} & \{A2\} & \{C1\} & \{C2\} & \{C3\} & \{B1,B2\} & \{B3,A3,C4\} \\
        \{A1\} & 0.0000 \\
        \{A2\} & 7.3485 & 0.0000 \\
        \{C1\} & 10.3441 & 4.1231 & 0.0000 \\
        \{C2\} & $\color{red}3.7417$ & 9.8995 & 11.8743 & 0.0000 \\
        \{C3\} & 5.4772 & 6.4807 & 9.6437 & 8.6023 & 0.0000 \\
        \{B1,B2\} & 3.7417 & 8.2462 & 10.8167 & 4.2426 & 8.8318 & 0.0000 \\
        \{B3,A3,C4\} & 4.5826 & 7.1414 & 7.8740 & 4.1231 & 6.4031 & 6.5574 & 0.0000 \\
    \end{matrix}
\end{align*}

Merge clusters \{A1\} and \{C2\}.
\begin{align*}
    \begin{matrix}
         & \{A2\} & \{C1\} & \{C3\} & \{B1,B2\} & \{B3,A3,C4\} & \{A1,C2\} \\
        \{A2\} & 0.0000 \\
        \{C1\} & 4.1231 & 0.0000 \\
        \{C3\} & 6.4807 & 9.6437 & 0.0000 \\
        \{B1,B2\} & 8.2462 & 10.8167 & 8.8318 & 0.0000 \\
        \{B3,A3,C4\} & 7.1414 & 7.8740 & 6.4031 & 6.5574 & 0.0000 \\
        \{A1,C2\} & 7.3485 & 10.3441 & 5.4772 & $\color{red}3.7417$ & 4.1231 & 0.0000 \\
    \end{matrix}
\end{align*}

Merge clusters \{B1,\ B2\} and \{A1,\ C2\}.
\begin{align*}
    \begin{matrix}
         & \{A2\} & \{C1\} & \{C3\} & \{B3,A3,C4\} & \{A1,B1,B2,C2\} \\
        \{A2\} & 0.0000 \\
        \{C1\} & $\color{red}4.1231$ & 0.0000 \\
        \{C3\} & 6.4807 & 9.6437 & 0.0000 \\
        \{B3,A3,C4\} & 7.1414 & 7.8740 & 6.4031 & 0.0000 \\
        \{A1,B1,B2,C2\} & 7.3485 & 10.3441 & 5.4772 & 4.1231 & 0.0000 \\
    \end{matrix}
\end{align*}

Merge clusters \{A2\} and \{C1\}.
\begin{align*}
    \begin{matrix}
         & \{C3\} & \{B3,A3,C4\} & \{A1,B1,B2,C2\} & \{A2,C1\} \\
        \{C3\} & 0.0000 \\
        \{B3,A3,C4\} & 6.4031 & 0.0000 \\
        \{A1,B1,B2,C2\} & 5.4772 & $\color{red}4.1231$ & 0.0000 \\
        \{A2,C1\} & 6.4807 & 7.1414 & 7.3485 & 0.0000 \\
    \end{matrix}
\end{align*}

Merge clusters \{B3, A3, C4\} and \{B1, B2, A1, C2\}.
\begin{align*}
    \begin{matrix}
         & \{C3\} & \{A2,C1\} & \{A1,A3,B1,B2,B3,C2,C4\} \\
        \{C3\} & 0.0000 \\
        \{A2,C1\} & 6.4807 & 0.0000 \\
        \{A1,A3,B1,B2,B3,C2,C4\} & $\color{red}5.4772$ & 7.1414 & 0.0000 \\
    \end{matrix}
\end{align*}

Merge clusters \{C3\} and \{A1, A3, B1, B2, B3, C2, C4\}.

\begin{align*}
    \begin{matrix}
         & \{A2,C1\} & \{A1,A3,B1,B2,B3,C2,C3,C4\} \\
        \{A2,C1\} & 0.0000 \\
        \{A1,A3,B1,B2,B3,C2,C3,C4\} & $\color{red}6.4807$ & 0.0000 \\
    \end{matrix}
\end{align*}

\textbf{The result of AGNES algorithm on the dataset generating 2 clusters is \{A2, C1\} and \{A1, A3, B1, B2, B3, C2, C3, C4\}.}





\subsection{Recommender System}
\textcolor{gray}{\it Table gives a User-Product rating matrix.}
\begin{table}[H]
    \centering
    \renewcommand{\arraystretch}{1.2}
    \setlength{\tabcolsep}{24pt}
    \color{gray}
    \caption{User-Product Rating Matrix}
    \begin{tabular}{ccccc}
        \bottomrule[1.5pt]
        \@ & Product 1 & Product 2 & Product 3 & Product 4 \\
        \midrule[0.75pt]
        User 1 & 1 & 1 & 5 & 3 \\
        User 2 & 3 & \textcolor{red}{?} & 5 & 4 \\
        User 3 & 1 & 3 & 1 & 1 \\
        User 4 & 4 & 3 & 2 & 1 \\
        User 5 & 2 & 2 & 2 & 4 \\
        \toprule[1.5pt]
    \end{tabular}
\end{table}



\subsubsection{Cosine similarity}
\textcolor{gray}{\it List the top 3 most similar users of user 2 based on Cosine Similarity.}

Define the Cosine Similarity as:
\[ \cos(\bm{x},\ \bm{y})=\sum_{i=1}^n\dfrac{x_iy_i}{|\bm{x}|\cdot|\bm{y}|} \]

\begin{align*}
    {\rm CosineSimilarity(User\ 2,\ User\ 1)}&=\dfrac{3\times1+5\times5+4\times3}{\sqrt{3^2+5^2+4^2}\cdot\sqrt{1^2+5^2+3^2}}\approx0.9561828874675149
    \\
    {\rm CosineSimilarity(User\ 2,\ User\ 3)}&\approx0.9797958971132713
    \\
    {\rm CosineSimilarity(User\ 2,\ User\ 4)}&\approx0.802377419802878
    \\
    {\rm CosineSimilarity(User\ 2,\ User\ 5)}&\approx0.9237604307034013
\end{align*}

The top 3 most similar users of User 2 are User 1, User 3 and User 5.



\subsubsection{User product rating}
\textcolor{gray}{\it Predict User 2's rating for Product 2.}

Assume that neighborhood size = 3.
\begin{align*}
    \overline{r}_{\rm User 2}&=\dfrac{3+5+4}{3}=4
    \\
    \overline{r}_{\rm User 1}&=\dfrac{1+1+5+3}{4}=2.5
    \\
    \overline{r}_{\rm User 3}&=\dfrac{1+3+1+1}{4}=1.5
    \\
    \overline{r}_{\rm User 5}&=\dfrac{2+2+2+4}{4}=2.5
\end{align*}
\begin{align*}
    r({\rm User 2,\ Product 2})&=\overline{r}_{\rm User 2}
    \\
    &\quad+\dfrac{\cos({\rm User 2,\ User 1})(r_{\rm User 1,\ Product 2}-\overline{r}_{\rm User 1})}{\cos({\rm User 2,\ User 1})+\cos({\rm User 2,\ User 3})+\cos({\rm User 2,\ User 5})}
    \\
    &\quad+\dfrac{\cos({\rm User 2,\ User 3})(r_{\rm User 3,\ Product 2}-\overline{r}_{\rm User 3})}{\cos({\rm User 2,\ User 1})+\cos({\rm User 2,\ User 3})+\cos({\rm User 2,\ User 5})}
    \\
    &\quad+\dfrac{\cos({\rm User 2,\ User 5})(r_{\rm User 5,\ Product 2}-\overline{r}_{\rm User 5})}{\cos({\rm User 2,\ User 1})+\cos({\rm User 2,\ User 3})+\cos({\rm User 2,\ User 5})}
\end{align*}
\begin{mathsize}
    \[ r({\rm User 2,\ Product 2})=4+\dfrac{0.9561828874675149\times(1-2.5)+0.9797958971132713\times(3-1.5)+0.9237604307034013\times(2-2.5)}{0.9561828874675149+0.9797958971132713+0.9237604307034013} \]
\end{mathsize}
\[ r({\rm User 2,\ Product 2})\approx3.850874269022923 \]





\section{Laboratory Assignment}
\subsection{Churn Management}
\textcolor{gray}{\it The goal of this assignment is to introduce churn management using decision trees, logistic regression and neural network. You will try different combinations of the parameters to see their impacts on the accuracy of your models for this specific data set. This data set contains summarized data records for each customer for a phone company. Our goal is to build a model so that this company can predict potential churners. Each row in ``churn\_training'' represents the customer record. The training data contains 2000 rows and the validation data contains 1033 records.}


\subsubsection{Decision tree classification}
\begin{figure}[H]
    \centering
    \includegraphics*[scale=0.5]{./figs/DecisionTree.png}
    \caption{Decision Tree}
\end{figure}



\subsubsection{Neural network}
\begin{figure}[H]
    \centering
    \includegraphics*[scale=0.5]{./figs/NeuralNetwork.png}
    \caption{Neural Network}
\end{figure}

\textcolor{red}{Note: due to the different initialization parameters of the neural network, the final results obtained may vary slightly.}



\subsubsection{Logistic regression}
\begin{figure}[H]
    \centering
    \includegraphics*[scale=0.5]{./figs/LogisticRegression.png}
    \caption{Logistic Regression}
\end{figure}


\subsubsection{Confusion matrices}


\begin{table}[H]
    \centering
    \renewcommand{\arraystretch}{1.2}
    \setlength{\tabcolsep}{12pt}
    \caption{Metrics}
    \begin{tabular}{ccccc}
        \bottomrule[1.5pt]
        \@ & Accuracy & Precision & Recall & F1 \\
        \midrule[0.75pt]
        Decision Tree & 0.94385 & 0.38806 & 0.60465 & 0.47273 \\
        Neural Network & 0.94773 & 0.40678 & 0.55814 & 0.47059 \\
        Logistic Regression & 0.93611 & 0.24444 & 0.25581 & 0.25000 \\
        \toprule[1.5pt]
    \end{tabular}
\end{table}

Because the ratio gap (approximately $900:30$) between positive samples and negative samples is huge, the accuracy is not a very good metric. We can consider three indicators: precision, recall and f1.
\begin{itemize}
    \item \textbf{Neural Network:} Shows strong performance, achieving the highest Accuracy (0.94773) and Precision (0.40678) among the three. This indicates it is the most reliable when it classifies a sample as positive.
    \item \textbf{Decision Tree:} Outperforms the other models in identifying positive instances, securing the highest Recall (0.60465). It also achieves the highest F1 Score (0.47273), marginally surpassing the Neural Network (0.47059), indicating a slightly better trade-off between precision and recall.
    \item \textbf{Logistic Regression:} Is the underperformer in this comparison. While its accuracy is comparable, it scores significantly lower in Precision, Recall, and F1 (0.25000), making it the least effective model for this specific task.
\end{itemize}

Conclusion: If the priority is minimizing false positives (Precision), the \textbf{Neural Network} is preferable. However, if the goal is to capture as many positive instances as possible (Recall) or to achieve the best overall balance (F1), the \textbf{Decision Tree} is the optimal choice.



\subsection{Market Basket Analysis}
\textcolor{gray}{\it Learn the use of market basket analysis for the purpose of making product purchase recommendations to the customers.}

\textcolor{gray}{\it The data set contains transactions from a large supermarket. Each transaction is made by someone holding the loyalty card. We limited the total number of categories in this supermarket data to 20 categories for simplicity. The field value for a certain product in the transaction basket is 1 if the customer has bought it and 0 if he/she has not. The file named ``Transactions'' has data for 46243 transactions.}

\textcolor{gray}{\it The data are available from the class web page. Your submission should consist only of those deliverables marked indicated by ``Hand-in''.}

\textcolor{gray}{\it Market basket analysis has the objective to discover individual products, or groups of products that tend to occur together in transactions. The knowledge obtained from a market basket analysis can be employed by a business to recognize products frequently sold together in order to determine recommendations and cross-sell and up-sell opportunities. It can also be used to improve the efficiency of a promotional campaign.}



\subsubsection{Association rules}
\textcolor{gray}{\it The list of association rules generated by the model. Sort the rules by lift, support, and confidence, respectively to see the rules identified. Hand-in: For each case, choose top 5 rules (note: make sure no redundant rules in the 5 rules) and give 2-3 lines comments. Many of the rules will be logically redundant and therefore will have to be eliminated after you think carefully about them.}

\begin{figure}[H]
	\centering
	\begin{minipage}{0.32\linewidth}
		\centering
		\includegraphics[width=1\linewidth]{./figs/Lift.png}
		\caption{Sorted by Lift}
	\end{minipage}
	\begin{minipage}{0.32\linewidth}
		\centering
		\includegraphics[width=1\linewidth]{./figs/Support.png}
		\caption{Sorted by Support}
	\end{minipage}
    \begin{minipage}{0.32\linewidth}
		\centering
		\includegraphics[width=1\linewidth]{./figs/Confidence.png}
		\caption{Sorted by Confidence}
	\end{minipage}
\end{figure}




\begin{lstlisting}[
    basicstyle = \ttfamily,
    columns = flexible,
    breaklines = true,
    xleftmargin = 2em,
    escapeinside = {!@}{@!}
]
// Lift
1. tomato source -> pasta
2. pasta + biscuits -> milk
3. pasta + water -> milk
4. juices -> milk
5. yoghurt -> milk

// Support
1. pasta -> milk
2. water -> milk
3. biscuits -> milk
4. brioches -> milk
5. yoghurt -> milk

// Confidence
1. biscuits + pasta -> milk
2. water + pasta -> milk
3. juices -> milk
4. tomato source -> pasta
5. yoghurt -> milk
\end{lstlisting}



In the Lift sorting, remove the rule of ``{\tt coffee + milk -> pasta}'' because in Westerners, coffee and milk are usually for breakfast, while pasta is for dinner, which is unreasonable.

In the Support sorting, ``{\tt $\varnothing$ -> milk}'' is removed because it is not an association rule; it merely indicates that milk is the best-selling item.

In the Confidence sorting, all five rules are relatively reasonable. For instance, shoppers who buy water, biscuits or yogurt and other beverages often purchase milk at the same time. Another example is that shoppers who buy tomato source are likely to purchase pasta.



\end{document}